\chapter*{不定积分每日一题}

\refer{《积分的方法与技巧》金玉明\ 顾新身\ 毛瑞庭\ /\ 编著}

\exercise{求积分$\int{\sec^{2n+2}{x}\d x}$}%在花括号内输入题目

\teacher{张晓鹏}%在花括号内输入你的名字

\postdate{2022年12月27日}%在花括号内输入你想要发布那一天的日期

\result{\sum_{k=0}^n\binom{n}{k}\frac{\tan^{2k+1}{x}}{2k+1}+C}%在花括号内输入本题结果

\analysis{
    又双叒叕是三角函数，我们前面探讨了$\sin{x}$和$\cos{x}$乘积的积分，今天我们来继续探索除法是怎么样的，最特殊的莫过于$\int{\sec^n{x}}$和$\int{\csc^n{x}}$了.今天先探讨一个相对简单的情况$\int{\sec^{2n+2}{x}\d x}$，如果一下子没有感觉的话，我们可以从简单的低次入手.
    \formula{\sec^2{x}\d x=\tan{x}+C}
    \formula{
    \begin{aligned}    
    \int{\sec^4{x}\d x}&=\int{\sec^2{x}\cdot\sec^2{x}\d x}=\int{\sec^2{x}\d\tan{x}}\\
    &=\int\qty(1+\tan^2{x}\d\tan{x})=\tan{x}+\frac{1}{3}\tan^3{x}+C
    \end{aligned}
    }
    欸，做了两个，熟练的小伙伴感觉就来了，如果是偶数的情况，我可以把$\sec^2{x}$放到$\d$后面，从而全部转化为$\tan{x}$积分，明白了这个，剩下的一切都不是问题.
    
    那我们直接推导了，转化为$\tan{x}$积分后二项展开，再分项积分即可.
    \formula{
        \begin{aligned}
            \int{\sec^{2n+2}{x}\d x}&=\int{\sec^{2n}{x}\d\tan{x}}=\int{\qty(1+\tan^2{x})^n\d\tan{x}}\\
            &=\int{\qty[1+\binom{n}{1}\tan^2{x}+\binom{n}{2}\tan^4{x}+\cdots+\binom{n}{n}\tan^{2n}{x}]\d\tan{x}}\\
            &=\tan{x}+\binom{n}{1}\frac{\tan^3{x}}{3}+\binom{n}{2}\frac{\tan^5{x}}{5}+\cdots+\binom{n}{n}\frac{\tan^{2n+1}{x}}{2n+1}+C\\
            &=\sum_{k=0}^n\binom{n}{k}\frac{\tan^{2k+1}{x}}{2k+1}+C
        \end{aligned}
    }
    跟之前都一样，我们不必直接记忆结论，会推导即可（如果你记忆力很猛，那也可以）.
}%在花括号内输入分析